%!TEX root = ../../exa-ma-deliverables.tex
% The June ANR review deck (slides.exa-ma.anr.20260611/content.typ and
% assets/exa-ma-ecosystem-links.svg) covers 2025 and early 2026. The rows
% below retain those dates and add evidence from 2026 where available.
\begingroup
\small
\setlength{\tabcolsep}{4pt}
\renewcommand{\arraystretch}{1.12}
\begin{longtable}{@{}
  >{\raggedright\arraybackslash}p{0.225\textwidth}
  >{\raggedright\arraybackslash}p{0.155\textwidth}
  >{\raggedright\arraybackslash}p{0.57\textwidth}@{}}
\caption{Exa-MA research collaborations and related projects, November 2026.}
\label{tab:exama-ecosystem-2026}\\
\toprule
\textbf{Project or institution} & \textbf{Status} & \textbf{Work and links to Exa-MA}\\
\midrule
\endfirsthead
\toprule
\textbf{Project or institution} & \textbf{Status} & \textbf{Work and links to Exa-MA}\\
\midrule
\endhead
\midrule
\multicolumn{3}{r@{}}{\emph{Continued on next page}}\\
\endfoot
\bottomrule
\endlastfoot

\multicolumn{3}{@{}l@{}}{\textbf{\href{https://numpex.org/projects/}{NumPEx targeted projects}}}\\*
PC2 --- Exa-SofT & Software collaboration &
WP3 works on composable linear algebra with Composyx. WP1 has
explored StarPU for FunTiDES. \href{https://exa-ma.numpex.fr/workpackages/wp7.html}{WP7
software integration}.\\[1pt]

PC3 --- Exa-DoST & Scientific collaboration &
Adaptive-mesh I/O and coupling Scimba with the GYSELA plasma
code. Exa-DoST also reports joint work on turbulence models. \href{https://numpex.org/wp-content/uploads/sites/6/2025/12/2025-11-05-01-NumPEx-and-Exa-DoST-Overview-Gabriel.pdf}{Exa-DoST
assembly overview}.\\[1pt]

PC4 --- Exa-AToW & Workflow collaboration &
The June ANR review records PC4's first deployment of Ktirio Urban
Building on Jean Zay.\\[1pt]

PC5 --- Exa-DI & Co-design &
Joint work on high-order discretisation, adaptive and unstructured
meshes, and application demonstrators. \href{https://numpex.org/exadi-development-and-integration/}{Exa-DI
programme and annual meeting}.\\

\addlinespace
\multicolumn{3}{@{}l@{}}{\textbf{\href{https://numpex.org/projects-call-2025/}{Projects funded through the 2025 NumPEx call}}}\\*
SAGE-HPC & Started in 2026 &
AI for HPC, focused on multifidelity optimisation. Inria and
Cemosis/IRMA share teams with Exa-MA. \href{https://numpex.org/wp-content/uploads/sites/6/2026/02/2026-01-21-1-SageHPC.pdf}{Assembly
slides}.\\[1pt]

DAiMOS & Started in 2026 &
HPC for AI, focused on distributed training. IRIT--Toulouse INP,
Inria Bordeaux, Sorbonne Université and Météo-France form the
consortium. \href{https://numpex.org/wp-content/uploads/sites/6/2026/02/2026-01-21-2-DAIMOS.pdf}{Assembly
slides}.\\[1pt]

KOKTAILS & Potential link &
GPU portability with Kokkos, including unstructured meshes. IFPEN
and Inria are involved.\\

\addlinespace
\multicolumn{3}{@{}l@{}}{\textbf{Research programmes and facilities}}\\*
HiDALGO2 / EuroHPC JU & Joint pilot and training &
Université de Strasbourg/Cemosis link WP1 and WP7 to the \href{https://www.hidalgo2.eu/urban-building-model/}{Ktirio Urban
Building pilot}. The teams and NCC France held a
\href{https://exa-ma.numpex.fr/training.html}{joint webinar on 12 January 2026}.\\[1pt]

JNL-G / LNCMI--Grenoble & Started in 2026 &
The LNCMI--IRMA digital twin reuses Exa-MA multiphysics and
model-reduction methods in Feel++ and HiFiMagnet.
\href{https://anr.fr/Project-ANR-25-CE46-1157}{ANR project summary}.\\[1pt]

% Funding shares from Exa-MA's 2025 ANR return, WP2 staffing.
\href{https://cluster-ia-enact.ai/}{ENACT AI cluster} & Co-funded WP2 PhD &
Exa-MA and the Grand Est AI cluster each fund half of Vincent
Italiano's doctoral contract.\\[1pt]

GENCI, AMD and CINES & Technical exchanges &
Exa-MA proposed three pilot topics with WP1, WP3 and WP5 at the
\href{https://exa-ma.numpex.fr/news/archive.html}{February workshop
in Saclay}.\\

\end{longtable}
\endgroup
